\documentclass[11pt]{article}
\usepackage{udec-pset}
\doc{4}{Potencial Eléctrico}

\begin{document}
\crearEncabezado

\begin{boxTeoria}{Potencial Eléctrico y Diferencia de Potencial}
    El \textbf{Potencial Eléctrico} ($V$) en un punto del espacio se define como la energía potencial eléctrica por unidad de carga. La \textbf{Diferencia de Potencial} ($\Delta V$), comúnmente llamada voltaje, representa el trabajo externo por unidad de carga necesario para mover una carga de prueba de un punto $A$ a un punto $B$ a velocidad constante:
    \begin{equation*}
        \Delta V = V_B - V_A = -\int_{A}^{B} \vec{E} \cdot d\vec{l} = \frac{W_{A \to B}}{q} \quad [\unit{\V}]
    \end{equation*}
    Para una carga puntual $q$ en el origen, el potencial a una distancia $r$ (asumiendo que $V \to 0$ cuando $r \to \infty$) es:
    \begin{equation*}
        V(r) = \frac{1}{4\pi\varepsilon_0} \frac{q}{r} \quad [\unit{\V}]
    \end{equation*}
    \begin{itemize}
        \item \textbf{Naturaleza Escalar:} A diferencia del campo eléctrico, el potencial es una magnitud escalar. Esto facilita su cálculo, ya que se evitan las descomposiciones vectoriales.
        \item \textbf{Superficies Equipotenciales:} Son regiones del espacio donde el potencial es constante. El vector de campo eléctrico $\vec{E}$ siempre es perpendicular a estas superficies, apuntando hacia donde el potencial disminuye.
    \end{itemize}

    El \textbf{Principio de Superposición} para el potencial dicta que el potencial neto en un punto, debido a una distribución de cargas discretas, es la suma \textit{algebraica} de los potenciales individuales: $V_{R} = V_{1} + V_{2} + \dots + V_{n}$.
\end{boxTeoria}

\begin{boxNota}{Unidades y Equivalencias Útiles}
    \begin{itemize}
        \item \textbf{El Voltio (V):} La unidad del SI para el potencial eléctrico. Equivale a un Joule por Coulomb ($\qty{1}{\V} = \qty{1}{\J\per\C}$).
        \item \textbf{Unidad equivalente para $\vec{E}$:} A partir de la definición de voltaje, el campo eléctrico suele expresarse comúnmente en Volts por metro (\unit{\V\per\meter}), lo cual es dimensionalmente idéntico a \unit{\N\per\C}.
        \item \textbf{Electrón-voltio (\unit{\electronvolt}):} Unidad de \textit{energía} (no de potencial) que equivale a la energía cinética adquirida por un electrón al moverse a través de una diferencia de potencial de \qty{1}{\V}. $\qty{1}{\electronvolt} = \qty{1.602e-19}{\J}$.
    \end{itemize}
\end{boxNota}

\problema
Una carga puntual $q$ es desplazada (por un agente externo) desde el origen $(0,0)$ hasta un punto $P(L, H)$ en una región con un campo eléctrico constante $\vec{E} = E_0 \hat{x}$.

Calcule
\begin{itemize}
    \item La diferencia de potencial $V_{OP}$.
    \item El trabajo realizado por el agente externo.
\end{itemize}

\problema
Dos cargas idénticas $+Q$ se ubican en el eje $y$, en las posiciones $y = d$ y $y = -d$. Encuentre el potencial en un punto $P(x, 0)$ sobre el eje $x$.

\problema
Un anillo de radio $R$ tiene una carga total $Q$ distribuida uniformemente ($\rho_l = \frac{Q}{2\pi R}$). Encuentre el potencial en un punto $P$ a una altura $z$ sobre su centro.


\problema
Considere dos esferas metálicas aisladas, de radios R y 3R. Si ambas esferas se encuentran al mismo potencial, ¿cuál
es la razón entre sus cargas? Si ambas esferas poseen la misma carga, ¿cuál es la razón entre sus potenciales?



\end{document}